\chapter{Lorentz Structures, Complex Units and Observational Spacetime}
\label{ch:ch02_lorentz}

The question of this chapter is which Lorentz algebra, which clock and which complex unit the algebra selects, and under which
hypotheses. A conversion constant ties the temporal triple to the spatial one; a clock is a unit vector of the temporal triple, and all
clocks are equivalent under rotations (\cref{rem:clockcov}). A clock selects an observable four-plane, and the restriction of physical
states to a single clock is an assumption, stated as \cref{assum:clockred} and analysed in \cref{ch:ch12_causal}. The joint commutant of
the interior Lorentz algebra and the internal bivectors fixes the imaginary unit up to sign (\cref{thm:i}), and the real Dirac theory
built on it is complex Dirac theory. Discrete symmetries and the state space close the chapter. Uniqueness statements hold under the
hypotheses they state; global covariance under a change of clock is not a local gauge symmetry.


% --- source: tex/s03_clock#0 ((head))


% --- source: tex/s03_clock#1 (Time, conversion, and the two square roots of the pairing)
\section{Time, conversion, and the two square roots of the pairing}
\label{sec:clock}

% --- source: tex/s03_clock#2 (One conversion constant)
\subsection{One conversion constant}

The split form is not dimensionally innocent. If $V^+$ carries times and $V^-$ lengths, then
$\eta(v,v)=\sum_i(t^i)^2-\sum_a(x^a)^2$ adds squared times to squared lengths, and a conversion is needed, a priori one
constant for each temporal axis.

\begin{proposition}[A pairing alone does not give isotropy]
\label{prop:KC}
A map exchanging $V^+$ and $V^-$ cannot carry a single dimension; it converts times to lengths one way and lengths to times
the other. In block form
\begin{equation}
\label{eq:KC}
\cK_C=\begin{pmatrix}0&C^{-1}\\ C&0\end{pmatrix},\qquad C:V^+\to V^-,\qquad [C]=LT^{-1},
\end{equation}
one has $\cK_C^2=\id$ for every invertible $C$, including $C=\operatorname{diag}(c_1,c_2,c_3)$ with distinct entries. An
involutive pairing therefore does not force $c_1=c_2=c_3$.
\end{proposition}

\begin{theorem}[One conversion constant]
\label{thm:onec}
If a common $SO(3)$ acts on $V^+$ and on $V^-$ by the vector representation and $C$ intertwines the two actions, then
$C=c\,\id$ with $[c]=LT^{-1}$.
\end{theorem}

\begin{proof}
The vector representation of $SO(3)$ is irreducible of real type, so its commutant is $\R$ and Schur's lemma places $C$ in
$\R\,\id$ \cite{GoodmanWallach}.
\end{proof}

% --- source: tex/s03_clock#3 (The two square roots of the pairing)
\subsection{The two square roots of the pairing}

Dimensional analysis fixes $[C]$ and leaves the sign of the lower-left block free. The two choices are different operators,
and the rest of the paper turns on the difference.

\begin{proposition}[The grading root and the binding root]
\label{prop:twoK}
In units with $c=1$, and in block form on $V^+\oplus V^-$, put
\[
\cK=\begin{pmatrix}0&-\id\\ -\id&0\end{pmatrix},\qquad
G=\begin{pmatrix}0&-\id\\ \id&0\end{pmatrix},
\]
that is $\cK e_i=-e_{i+3}$, $\cK e_{i+3}=-e_i$, and $Ge_i=e_{i+3}$, $Ge_{i+3}=-e_i$. Then:
\begin{center}
\begin{tabular}{@{}lcccl@{}}
\toprule
& square & $\eta$-adjoint & $\eta(X\cdot,X\cdot)$ & role\\
\midrule
$\cK$ & $+\id$ & $\cK^\dagger=-\cK$ & $-\eta$ & an element of $\so{3,3}$; grades\\
$G$ & $-\id$ & $G^\dagger=+G$ & $-\eta$ & a complex structure outside $\so{3,3}$; binds\\
\bottomrule
\end{tabular}
\end{center}
with $X^\dagger:=\eta^{-1}X^{\!\top}\eta$. Both are anti-isometries. The eigenspaces of $\cK$ are the two totally null three-planes
spanned by $e_i\mp e_{i+3}$.
\end{proposition}

Both roots can be built inside the algebra, and they are related by a single trivector.

\begin{lemma}[Both roots from Clifford data]
\label{lem:Gclifford}
Let $\varpi=\sum_i e_ie_{i+3}$ and $T=e_1e_2e_3$, and define $\cK(v)=\tfrac12[\varpi,v]$ and $S(v)=TvT^{-1}$ on vectors. Then
$\cK(e_i)=-e_{i+3}$, $\cK(e_{i+3})=-e_i$, so $\cK^2=\id$; $S$ fixes the temporal and negates the spatial generators, so $S=\eta$ as a
map of $V$; and
\[
G=-\cK\circ S .
\]
\end{lemma}

\begin{proof}
$[e_ie_{i+3},e_i]=e_ie_{i+3}e_i-e_i^2e_{i+3}=-2e_{i+3}$ and $[e_ie_{i+3},e_{i+3}]=e_ie_{i+3}^2-e_{i+3}e_ie_{i+3}=-2e_i$, while the
other pole operators commute with both by \eqref{eq:bladecomm}. $T$ is a product of the three temporal generators, so it
commutes with each of them up to the sign $(-1)^{3-1}=+1$ and anticommutes with each spatial generator. Finally
$-\cK(S e_i)=-\cK(e_i)=e_{i+3}=Ge_i$ and $-\cK(Se_{i+3})=\cK(e_{i+3})=-e_i=Ge_{i+3}$.
\end{proof}

So the root that grades is $\tfrac12\ad_{\varpi}$, the adjoint action of the sum of the pole operators, and the root that binds
is the same map composed with the Cartan involution $S=\Ad_{e_1e_2e_3}$. That the two roots are distinguished exactly by the
Cartan involution is the cleanest single statement the flat and the curved analyses below jointly support.

\paragraph{Shift operators.} A clock coordinate $\lambda$ has dimension $T$, the operator $G$ is dimensionless, and their sum
requires a scale. With a length $\ell$,
\begin{equation}
\label{eq:kappa}
\kappa=\frac c\ell,\qquad \alpha_\pm=\partial_\lambda\pm\kappa G .
\end{equation}
The pointwise $\eta$-adjoint of $\kappa G$ is $\kappa G$, while the formal adjoint of the differential operator, after an integration by
parts with compatible boundary conditions, is $(\partial_\lambda+\kappa G)^*=-\partial_\lambda+\kappa G^\dagger=-\alpha_-$. The two
adjoints must be kept apart: the factorisations below are insensitive to the overall sign, positivity statements are not.

% --- source: tex/s03_clock#4 (The clock)
\subsection{The clock}
\label{sec:clocksimplex}

A clock is a unit vector of $V^+$. Which one? Let $w=(w_1,w_2,w_3)$ be normalised non-negative weights on the three temporal
axes, so that $w$ ranges over the two-simplex $\Delta^2$. The map $x_i=2\sqrt{w_i}$ sends $\Delta^2$ onto the part of the sphere
$\sum_ix_i^2=4$ in the closed positive octant, and since $dx_i=dw_i/\sqrt{w_i}$ the Euclidean metric pulls back to
\[
\sum_idx_i^2=\sum_i\frac{dw_i^2}{w_i},
\]
the Fisher--Rao metric. \emph{Given} that a clock is described by probability weights on the three axes, the metric is not a further
choice: by \v{C}encov's theorem it is, up to scale, the only metric on a finite simplex invariant under Markov morphisms
\cite{Cencov,AmariNagaoka}. That description is itself a modelling choice, not a consequence of the algebra. It restricts the clocks to
the closed positive octant of a round sphere of radius two, and the octant is preserved by the permutations of the axes but not by the
rotations of $V^+$, under which every unit vector of $V^+$ is equivalent to every other (\cref{rem:clockcov}).

\begin{proposition}[The symmetric clock]
\label{prop:clock}
The barycentre $w=(\tfrac13,\tfrac13,\tfrac13)$ is the unique point of $\Delta^2$ fixed by the permutations of the three axes, and the
unique maximum of the entropy $-\sum_iw_i\log w_i$, with value $\log3$ and Hessian eigenvalues $-9$ and $-3$ in the chart
$(w_1,w_2)$. Its amplitude direction is the unit vector
\[
\widehat W=\frac{e_1+e_2+e_3}{\sqrt3},\qquad \widehat W^2=\id,\qquad
\lambda:=\eta(\widehat W,v)=\frac{t^1+t^2+t^3}{\sqrt3}=\sqrt3\,\bar t .
\]
\end{proposition}

The permutation argument carries the weight here; the entropy argument is Jaynes' principle and is a hypothesis, recorded as
such. What the proposition says is modest: if a clock is to treat the three temporal axes alike, it is $\widehat W$, and the
physical time it reads is $\sqrt3$ times the mean of the three constituent times. It does not say that $\lambda$ is a proper time,
that it is invariant under a group, or that the directions orthogonal to $\widehat W$ are unphysical. Those directions,
$\widehat W^\perp\cap V^+=\{t^1+t^2+t^3=0\}$, spanned by $e_1-e_2$ and $e_1+e_2-2e_3$, are the \emph{relative times}.

\begin{remark}[Clocks: covariance, not selection]
\label{rem:clockcov}
Five objects must be kept apart: a unit clock $\hat n\in V^+$ ($\hat n^2=\id$), the unnormalised sum $e_1+e_2+e_3$ (square $3\,\id$), the temporal trivector
$T=e_1e_2e_3$ ($T^2=-\id$, not a vector), the discrete permutations of the axes, and the continuous rotations $SO(3)$ of $V^+$. The rotor
$R=(\id+\widehat We_1)/|\id+\widehat We_1|$ of $\mathrm{Spin}(V^+)$ carries the axial clock $e_1$ to $\widehat W$; it commutes with $T$ and carries
$J(e_1)=Te_1$ to $J(\widehat W)=T\widehat W$ and the free Dirac Hamiltonian $H_{e_1}(\bp)=me_1+\sum_ap_ae_1e_{a+3}$ to $H_{\widehat W}(\bp)$. For every unit
clock $J(\hat n)=T\hat n$, $J(\hat n)^2=-\id$ and $\hat nJ(\hat n)=T$. So the axial and the symmetric clock lie on one $SO(3)$ orbit, although not on one orbit of the
permutations, under which an odd permutation reverses $T$ and with it $J(\hat n)$. Three cautions follow. (a) This is a \emph{global} covariance of
the free data: it shows that physically equivalent descriptions exist, not that a change of clock is a gauge redundancy, which would
need an action, a redundancy of its solutions and the invariance of the observables. A position-dependent rotor $R(x)$ adds the term
$R^{-1}\partial_\mu R\in\so3_+$ to the Dirac operator, so a local clock symmetry requires a connection with its transformation law, which is
not constructed here. (b) Structures that are \emph{transported} with the rotor keep their algebraic properties; structures held
\emph{fixed} need not be invariant. The binding root $G$ is preserved by a rotation acting alike on the temporal and the spatial triple and
is changed by a rotation of the temporal triple alone; a matter projector chosen once, such as $P_*$ of \cref{sec:pstar}, is not
invariant under clock rotations unless it is transported. (c) Neither the permutation symmetry nor the entropy maximum selects a clock
dynamically; they single out $\widehat W$ only within the chosen octant description. Global covariance establishes neither a local
clock gauge symmetry, nor a dynamical clock selection, nor a well-posed interacting theory with three times (\cref{sec:baseline}).
\end{remark}

\begin{figure}[t]
\centering
\begin{tikzpicture}[scale=1.15]
% simplex
\begin{scope}
\coordinate (A) at (0,0); \coordinate (B) at (3,0); \coordinate (Cc) at (1.5,2.6);
\fill[tcol!8] (A)--(B)--(Cc)--cycle; \draw[tcol,thick] (A)--(B)--(Cc)--cycle;
\node[below left] at (A) {$e_1$}; \node[below right] at (B) {$e_2$}; \node[above] at (Cc) {$e_3$};
\coordinate (bar) at (1.5,0.8667);
\fill[scol] (bar) circle (2.3pt);
\node at (1.5,0.2) {\footnotesize $w=(\tfrac13,\tfrac13,\tfrac13)$};
\foreach \r in {0.22,0.42}{\draw[gray!60,thin] (bar) circle (\r);}
\node at (1.5,-0.6) {\small simplex $\Delta^2$ of clock weights};
\end{scope}
\draw[->,thick] (3.5,1.2) -- node[above] {\small $x_i=2\sqrt{w_i}$} (5.3,1.2);
% octant of sphere, drawn in an oblique projection
\begin{scope}[shift={(7.3,0)}]
\coordinate (O) at (0,0);
\coordinate (X) at (2.6,-0.25); \coordinate (Y) at (-0.9,-0.9); \coordinate (Z) at (0,2.6);
\draw[->,gray] (O)--($1.25*(X)$) node[right] {\small $x_2$};
\draw[->,gray] (O)--($1.35*(Y)$) node[left] {\small $x_1$};
\draw[->,gray] (O)--($1.15*(Z)$) node[above] {\small $x_3$};
\fill[tcol!12] (X) .. controls ($(X)+(0,1.8)$) and ($(Z)+(1.2,0)$) .. (Z)
  .. controls ($(Z)+(-0.8,-0.6)$) and ($(Y)+(-0.5,1.3)$) .. (Y)
  .. controls ($(Y)+(1.1,-0.2)$) and ($(X)+(-1.2,-0.6)$) .. (X);
\draw[tcol,thick] (X) .. controls ($(X)+(0,1.8)$) and ($(Z)+(1.2,0)$) .. (Z);
\draw[tcol,thick] (Z) .. controls ($(Z)+(-0.8,-0.6)$) and ($(Y)+(-0.5,1.3)$) .. (Y);
\draw[tcol,thick] (Y) .. controls ($(Y)+(1.1,-0.2)$) and ($(X)+(-1.2,-0.6)$) .. (X);
\coordinate (Wp) at (0.72,0.95);
\fill[scol] (Wp) circle (2.3pt);
\draw[->,scol,thick] (O) -- (Wp) node[above right] {\small $\widehat W$};
\node[align=center,anchor=north] at (1.4,-1.45) {\small octant of the sphere $\sum x_i^2=4$\\[-1pt] \small (Fisher--Rao metric)};
\end{scope}
\end{tikzpicture}
\caption{The space of clocks. Normalised weights on the three temporal axes form a simplex; the amplitude map carries it onto an
octant of a round sphere, whose round metric is the Fisher--Rao metric. The barycentre, the only permutation-invariant point and
the entropy maximum, corresponds to the symmetric clock $\widehat W=(e_1+e_2+e_3)/\sqrt3$ (\cref{prop:clock}).}
\label{fig:clock}
\end{figure}

% --- source: tex/s03_clock#5 (The spectrum of a conversion operator)
\subsection{The spectrum of a conversion operator}

An operator that converts times into lengths is a mixed bivector $B_C=\sum_{i,a}C_{ia}\,e_ie_{a+3}$, symmetric by
\cref{thm:sig}. Its square controls the factorised interior dynamics of \cref{sec:shifted}, where the equation of motion becomes
$\ddot X=\kappa^2B_C^2X$, so its spectrum is worth knowing exactly.

\begin{theorem}[Spectrum of a conversion operator]
\label{thm:spectrumBC}
Write $C=\sum_k\sigma_ku_kv_k^{\!\top}$ in singular-value form and put $U_k=\sum_iu_k^ie_i$, $V_k=\sum_av_k^ae_{a+3}$ and
$B_k=U_kV_k$. Then $B_k^2=\id$, the $B_k$ commute, and
\[
B_C^2=\|C\|_F^2\,\id+2\sum_{k<l}\sigma_k\sigma_l\,\Omega_{kl},\qquad \Omega_{kl}=B_kB_l,\qquad \Omega_{12}\Omega_{13}=\Omega_{23}.
\]
Consequently
\[
\spec\bigl(B_C^2\bigr)=\bigl\{(\eta_1\sigma_1+\eta_2\sigma_2+\eta_3\sigma_3)^2:\ \eta\in\{\pm1\}^3\bigr\},
\]
four values, each of multiplicity two on $\R^8$.
\end{theorem}

\begin{proof}
$U_k$ and $V_k$ are orthogonal unit vectors of square $+\id$ and $-\id$, so $B_k^2=-U_k^2V_k^2=\id$; for $k\neq l$ the four vectors
$U_k,V_k,U_l,V_l$ are mutually orthogonal and $B_k,B_l$ commute by \eqref{eq:bladecomm}. Expanding the square gives the formula. The
seven nontrivial products of the commuting involutions $B_k$ are, in a rotated orthonormal basis of $V$, blades of grade two, four
or six, hence traceless; the joint eigenspace with signs $\eta$ therefore has dimension $\tfrac18\Tr\prod_k(\id+\eta_kB_k)=1$, and on it
$B_C=\eta\cdot\sigma$.
\end{proof}

\begin{corollary}[Rank counts clocks]
\label{cor:rankclock}
$B_C^2$ is a multiple of $\id$ if and only if $\rank C\le1$, and some eigenvalue of $B_C^2$ vanishes if and only if one singular value
equals the sum of the other two. The rank of $C$ is the number of independent temporal directions the conversion uses.
\end{corollary}

\begin{example}[Conversion along one clock is dynamically blind]
\label{ex:blindclock}
For any unit clock $n\in V^+$ the operators $G_a(n):=n\,e_{a+3}$, which convert along $n$, have $C=n\,\hat e_a^{\!\top}$ of rank one and singular
values $(1,0,0)$. Hence $G_a(n)^2=\id$, the factorised equation of motion is $\ddot X=+\kappa^2X$, whose solutions grow or decay
exponentially and do not oscillate, and the conversion operator does not appear in it at all. This holds for every clock, the axial
$e_1$ and the symmetric $\widehat W$ alike, so it does not single out $\widehat W$: the blindness is a property of converting along a single
temporal direction, forced by the rank. A conversion visible in the equations of motion needs at least two independent temporal directions
in $C$, that is, a relative time with respect to the clock used. Raising the rank does not
produce oscillation either: the eigenvalues $(\eta\cdot\sigma)^2$ are real and non-negative at every rank, so higher rank makes the
instability anisotropic, with four rates $\kappa|\eta\cdot\sigma|$ instead of one, and can create flat directions, but never a bounded
oscillation. Bounded motion has to come from elsewhere, and \cref{sec:neutral} shows where.
\end{example}

\section{The observational four-plane of a clock}
\label{sec:w:theobservationalfourplan}

This section fixes a clock and the four-plane it selects. That physical states and observables use a single clock is \cref{assum:clockred}, stated at the end of the section; it is an assumption, analysed in \cref{ch:ch12_causal}.

% --- source: two_readings/r2_fourmomentum#2 (Configuration, clock and the selected four-plane)
\subsubsection{Configuration, clock and the selected four-plane}
\label{r2:sec:clock}

\begin{definition}[Configuration and clock]
\label{r2:def:config}
The classical configuration is $Y=\sum_ic\,t^ie_i+\sum_ax^as_a$, three converted times on $V^+$ and three lengths on $V^-$, with
$\eta(Y,Y)=c^2|t|^2-|x|^2$ in $\mathrm m^2$. A \emph{clock} is a unit vector $\hat n\in V^+$. It defines
\[
E_{\hat n}=\R \hat n\oplus V^-,\qquad U_{\hat n}=\hat n^\perp\cap V^+,\qquad V=E_{\hat n}\oplus U_{\hat n} ,
\]
an oriented orthonormal frame $(\hat n,u_1,u_2)$ of $V^+$, the time $t_{\hat n}:=\eta(\hat n,Y)/c=\hat n\cdot t$ read by the clock and the surplus coordinates
$u^a:=\eta(u_a,Y)$, in $\mathrm m$. With $\pi_\pm$ the projectors onto $V^\pm$ put $\pr_{\hat n}=\hat n\,\eta(\hat n,\cdot)+\pi_-$ and $\pr^{\perp}_{\hat n}=\pi_+-\hat n\,\eta(\hat n,\cdot)$.
The clocks of non-negative weight form the Fisher--Rao octant $\hat n_i=\sqrt{w_i}$, $w\in\Delta^2$, which is the octant of radius two of \cref{sec:clocksimplex} scaled by one half; the symmetric clock is $\widehat W=(e_1+e_2+e_3)/\sqrt3$ (\cref{prop:clock}).
\end{definition}

This is constructed: a reading rule. The splitting $V^+\oplus V^-$ is
background structure, not invariant under $SO(3,3)$.

\begin{proposition}[The four-planes that contain the three lengths]
\label{r2:prop:planes}
\begin{enumerate}[label=(\roman*),leftmargin=2.2em]
\item $\pr_{\hat n}$ and $\pr^{\perp}_{\hat n}$ are complementary, $\eta$-self-adjoint, $\eta$-orthogonal projectors onto $E_{\hat n}$ and $U_{\hat n}$; $\eta|_{E_{\hat n}}$ has
signature $(1,3)$ and $\eta|_{U_{\hat n}}$ is positive definite.
\item Every four-plane $E\subset V$ that contains $V^-$ is $E_{\hat n}$ for exactly one pair $\pm \hat n$. Hence the four-planes containing the three lengths are
parametrised by the clock sphere modulo $\pm$, a two-dimensional family inside the eight-dimensional space of Lorentzian four-planes of
\cref{thm:selector}(iii).
\end{enumerate}
\end{proposition}

\begin{proof}
(i) $\pr_{\hat n}+\pr^{\perp}_{\hat n}=\pi_++\pi_-=\id$, $\pr_{\hat n}n=\hat n$, $\pr_{\hat n}s_a=s_a$, $\pr^{\perp}_{\hat n}u_a=u_a$, and the Gram matrices are
$\operatorname{diag}(1,-1,-1,-1)$ on $(\hat n,s_1,s_2,s_3)$ and $\id$ on $(u_1,u_2)$. (ii) Since $V^-$ is nondegenerate, $E=V^-\oplus(E\cap V^+)$, and
$E\cap V^+$ is a line on which $\eta$ is positive, spanned by a unit $\hat n$ determined up to sign.
\end{proof}

Fixed (r2-1a, symbolic for a general frame in the quaternion parametrisation of $SO(3)$; r2-1b). Taking $X$ ``in the three directions $(-,-,-)$''
therefore leaves exactly one choice, that of a clock, and every observer of the three lengths is an $E_{\hat n}$.

\paragraph{The six-momentum.} The covector $K=K_A\,dY^A$ conjugate to $Y$ has the raised form
\begin{equation}
\label{r2:eq:K}
K^\sharp=\eta^{-1}K=\frac Ecn+\sum_ap_as_a+\sum_aq_au_a,
\end{equation}
so that $E/c=\eta(\hat n,K^\sharp)$, $p_a=-\eta(s_a,K^\sharp)$ and $q_a=\eta(u_a,K^\sharp)$.
With the phase-space form $K_A\,dY^A=E\,dt_{\hat n}-\sum_ap_a\,dx^a+\sum_aq_a\,du^a$ and $\{Y^A,K_B\}=\delta^A_B$ one finds
\[
\{ct_{\hat n},E/c\}=1,\qquad \{x^a,p_b\}=-\delta_{ab},\qquad \{u^a,q_b\}=\delta_{ab},
\]
and all other brackets among these twelve functions vanish (fixed; r2-5a, for the clocks $e_1$, $(2,3,6)/7$ and $(-1,4,8)/9$). So $\bp$ is conjugate to the three lengths, $E/c$
to the converted time $ct_{\hat n}$ of the clock and $q_a$ to the surplus coordinates; the relative sign between the first two is the Lorentzian lowering
$K_{a+3}=-p_a$. This is the convention of \cref{thm:wlgauge}; the action $-Mc\int ds$ has the opposite overall
sign, which exchanges which pair carries the minus sign and changes no equation of motion.

\begin{proposition}[Covariance of the selected projection]
\label{r2:prop:cov}
\begin{enumerate}[label=(\roman*),leftmargin=2.2em]
\item An element of $O(V,\eta)$ commutes with $\pr_{\hat n}$ iff it preserves $E_{\hat n}$ and $U_{\hat n}$; the commutant of $\pr_{\hat n}$ in $\so{3,3}$ is
$\lint(\hat n)\oplus\so{U_{\hat n}}$, of dimension $7$, where $\lint(\hat n)$ is spanned by the six bivectors of $E_{\hat n}$.
\item For $\Lambda\in SO^+(E_{\hat n})$ extended by the identity on $U_{\hat n}$, $\pr_{\hat n}(\Lambda K^\sharp)=\Lambda\,\pr_{\hat n}K^\sharp$ and
$q_a(\Lambda K^\sharp)=q_a(K^\sharp)$: the projection is a four-vector and the $q_a$ are Lorentz scalars.
\item The rotations $R_\theta=\exp(\theta\,u_1\wedge u_2)$ of $U_{\hat n}$ commute with $\pr_{\hat n}$ and turn $(u^1,u^2)$ and $(q_1,q_2)$ through the same
angle; their cotangent lift preserves $K_A\,dY^A$.
\item The eight generators of $E_{\hat n}\wedge U_{\hat n}$ do not commute with $\pr_{\hat n}$, and their eight commutators are linearly independent. For instance
at $\hat n=e_1$ the rotation $\exp(\varphi\,e_1\wedge e_2)$ carries the rest momentum $K^\sharp=(E/c)\,e_1$ to one with $|q_1|=(E/c)|\sin\varphi|$.
\end{enumerate}
\end{proposition}

\begin{proof}
(i) A map commuting with a projector preserves its image and kernel; an isometry preserving $E_{\hat n}$ preserves $E_{\hat n}^\perp=U_{\hat n}$. On Lie algebras this is
$\so{E_{\hat n}}\oplus\so{U_{\hat n}}$, of dimension $6+1$. (ii) $\Lambda$ commutes with $\pr_{\hat n}$ by (i), and $\eta(u_a,\Lambda K^\sharp)=\eta(\Lambda^{-1}u_a,K^\sharp)
=\eta(u_a,K^\sharp)$. (iii) $\eta$ is positive on $U_{\hat n}$, so contravariant and covariant components of $U_{\hat n}$ rotate alike; for an isometry $R$ the
contragredient on $K_A$ is $R^{-\top}=\eta R\eta^{-1}$, so $K^\sharp$ moves as $Y$ and $K_A\,dY^A$ is invariant. (iv) A generator of $E_{\hat n}\wedge U_{\hat n}$
maps $U_{\hat n}$ into $E_{\hat n}$; the example is a direct computation.
\end{proof}

Nothing here gives covariance under the whole
of $SO(3,3)$ (r2-2a for seven clocks including $\widehat W$, the dimension count for six rational clocks; r2-2b for a finite boost and the $SO(2)$ at
$\hat n=(2,3,6)/7$; r2-2c). Part (i) is the case $C=J(\hat n)$ of \cref{thm:selector}(iii).

\begin{proposition}[The clock, the surplus plane and $J(\hat n)$]
\label{r2:prop:J}
Let $I_T=e_1e_2e_3$.
\begin{enumerate}[label=(\roman*),leftmargin=2.2em]
\item $I_T^2=-\id$; $I_T$ commutes with $e_1,e_2,e_3$ and anticommutes with $e_4,e_5,e_6$.
\item $J(\hat n):=I_T\hat n=\hat n_1e_2e_3+\hat n_2e_3e_1+\hat n_3e_1e_2$ satisfies $J(\hat n)^2=-|\hat n|^2\id$ and commutes with $\hat n$ and with $s_1,s_2,s_3$, hence with all of $\lint(\hat n)$.
\item For every oriented orthonormal frame $(\hat n,u_1,u_2)$ of $V^+$, $nu_1u_2=I_T$; hence $J(\hat n)=u_1u_2$ is the oriented unit bivector of $U_{\hat n}$, and
$J(e_1)=e_2e_3=J$.
\item Under time reversal $\hat n\mapsto-\hat n$: $\pr_{-\hat n}=\pr_{\hat n}$, $\pr^{\perp}_{-\hat n}=\pr^{\perp}_{\hat n}$, $E(-\hat n)=-E(\hat n)$ for the energy
$E(\hat n):=c\,\eta(\hat n,K^\sharp)$ read by the clock, $J(-\hat n)=-J(\hat n)$ and
$\Omega_{-\hat n}=-\Omega_{\hat n}$ for $\Omega_{\hat n}:=ns_1s_2s_3$, so that $\Gamma_{\hat n}=J(\hat n)\Omega_{\hat n}$ is unchanged.
\end{enumerate}
\end{proposition}

\begin{proof}
(i) is the blade rule (\cref{lem:blade}): $|I||J|-|I\cap J|$ equals $2$ for $e_i$ and $3$ for $s_a$. (ii) $I_T$ and $\hat n$
commute, so $J(\hat n)^2=I_T^2\hat n^2=-|\hat n|^2$; $I_T$ and $\hat n$ both commute with $\hat n$ and both anticommute with each $s_a$. (iii) $nu_1u_2$ is the volume element of
an oriented orthonormal frame of $V^+$, which is $\det(\hat n,u_1,u_2)\,e_1e_2e_3=I_T$; then $u_1u_2=\hat n^{-1}I_T=\hat nI_T=I_T\hat n$. (iv) $\pr_{\hat n}$ is even in $\hat n$,
$J$ and $\Omega$ are odd.
\end{proof}

Checks: r2-3a and r2-3c for symbolic $\hat n$, r2-3b for a
symbolic frame, r2-3d). The formula $J(\hat n)=I_T\hat n$ is the identity $J=\gamma^0T$ of \cref{rem:ambient} with $\gamma^0=\hat n$. Item (iv) says that time orientation and the sign of the complex structure are tracked together.

\begin{theorem}[The selected complex structure at every clock]
\label{r2:thm:i}
Let $\gamma^0=\hat n$, $\gamma^a=s_a$ and $\Omega_{\hat n}=ns_1s_2s_3$.
\begin{enumerate}[label=(\roman*),leftmargin=2.2em]
\item The joint commutant of $\lint(\hat n)$ and of the surplus bivector $J(\hat n)$ in $\End_\R(S)=M_8(\R)$ is $\operatorname{span}\{\id,\Omega_{\hat n},J(\hat n),\Omega_{\hat n}J(\hat n)\}
\cong\C\oplus\C$. Its solutions of $x^2=-\id$ are $\pm\Omega_{\hat n},\pm J(\hat n)$, and only $\pm J(\hat n)$ commute with the four $\gamma^\mu$.
\item The commutant of the four $\gamma^\mu$ alone is $\operatorname{span}\{\id,J(\hat n),u_1\Omega_{\hat n},u_2\Omega_{\hat n}\}\cong\HH$, and the elements of it that
commute with $I_T$ are $\operatorname{span}\{\id,J(\hat n)\}$.
\item The mass form $\gamma^0J(\hat n)=\hat nI_T\hat n=I_T$ is the same for every clock; the kinetic kernels are $I_TJ(\hat n)\gamma^\mu=-\hat n\gamma^\mu$; and
$H_{\hat n}(\bp)=c\,ns_ap_a+mc^2\hat n$ is symmetric, commutes with $J(\hat n)$ and satisfies $H_{\hat n}(\bp)^2=(c^2|\bp|^2+m^2c^4)\id$.
\end{enumerate}
\end{theorem}

\begin{proof}
A rotor $R$ of $\mathrm{Spin}(V^+)$ with $R\hat nR^{-1}$ prescribed exists for every unit $\hat n$, and since temporal bivectors commute with each $s_a$ by the
blade rule, $\Ad_R$ fixes $s_1,s_2,s_3$ and $I_T$. It therefore carries $(e_1,e_2,e_3)$ to an oriented frame $(\hat n,u_1,u_2)$, $\lint(e_1)$ to $\lint(\hat n)$,
$e_2e_3$ to $u_1u_2=J(\hat n)$ and $\Omega$ to $\Omega_{\hat n}$. Conjugation is an automorphism of $M_8(\R)$, so (i) and (ii) are \cref{thm:i} and \cref{rem:ambient} transported. For (ii), $I_T$ commutes with $u_a$ and anticommutes with $\Omega_{\hat n}$. (iii) $\hat n$ commutes with $I_T$ and $\hat n^2=\id$; the square of
$H_{\hat n}$ is that of equation~\cref{eq:Hp} with $\gamma^0=\hat n$.
\end{proof}

Fixed. Each part is checked directly, not only through the transport: r2-4a,b compute the commutants and solve $x^2=-\id$ exactly at
$\hat n=e_1,e_2,(1,2,2)/3,(2,3,6)/7,(-1,4,8)/9$; r2-4c is symbolic in a general frame; r2-4d checks the transport by the rotor
$\cos\frac\theta2-\sin\frac\theta2\,e_1e_2$. So the Dirac spinor read by the clock $\hat n$ has imaginary unit $J(\hat n)$, its sign fixed by the orientation of
$V^+$ and of $\hat n$, and its energy for momentum $\bp$ is the $E$ of the four-momentum below. The commutant conditions therefore hold at every clock $\hat n$. That the mass form does not depend on the clock is consistent with \cref{cor:fullmass}, where
$T=e_1e_2e_3$ is the mass kernel invariant under all of $\so{3,3}$.

\begin{remark}[Three readings of $(-,-,-)$]
\label{r2:rem:minus}
The sign triple has three meanings in the arena, and only the first bears on the four-momentum. (a) The values of $\eta$ on $V^-$, the reading used
here (interpretive: a reading of the question). (b) The sign sector $(-,-,-)$ of the temporal $\R^3$, which was proposed, in an earlier and superseded version of the programme, as the carrier of the Lorentzian interior: by
\cref{r2:prop:J}(iv) the clocks of that sector give the same four-planes and surplus planes as those of the Fisher--Rao octant $(+,+,+)$, with the
opposite time orientation, the opposite energy and the opposite complex structure, so the sector carries nothing the Fisher--Rao octant does not
(fixed). (c) In the spinor module the joint eigenline of the pole operators with signs $(-,-,-)$ is the filled Fock line $N=3$ (r2-0b), a spinor and not a
direction of $V$. The octant readings are the subject of \cref{r3:sec:main}.
\end{remark}


\subsection{The clock reduction is an assumption}
\label{sec:clockreduction}

The observable four-plane $E_{\hat n}$ of a clock and the four-momentum it reads are used throughout the following chapters. They rest on one
step that the algebra does not supply, and it is recorded here once.

\begin{assumption}[Clock reduction]
\label{assum:clockred}
Physical states and observables are those of the reduced shell of a single clock $\hat n$: the temporal part of every six-momentum is
parallel to $\hat n$, so that momenta lie in $E_{\hat n}=\R\hat n\oplus V^-$, of signature $(1,3)$, and the surplus temporal momenta vanish.
\end{assumption}

The assumption is imposed, not derived. \Cref{sec:constraints} shows where the constraints that implement it can come from, a worldline
gauge symmetry chosen rather than derived (\cref{thm:wlgauge}); that analysis establishes constraint preservation for a worldline in a
background field only, and the constraints remove every surplus-momentum mode, the stable ones included. Statements below that use the
four-momentum of a clock, the selected four-plane or the absence of uphill decays are statements under \cref{assum:clockred}, and they cite
it.

% --- source: tex/s04_dirac#0 ((head))


% --- source: tex/s04_dirac#1 (The imaginary unit and the Dirac structure)
\section{The imaginary unit and the Dirac structure}
\label{sec:dirac}

% --- source: tex/s04_dirac#2 (The complex structure is selected)
\subsection{The complex structure is selected}

At $n$ constituents write $\Omega=\gamma^0\gamma^1\gamma^2\gamma^3=e^{(1)}_1e^{(1)}_4e^{(1)}_5e^{(1)}_6$ for the volume element of the
interior four-plane and $\wI$ for the product, in increasing order, of the remaining $6n-4$ generators, the \emph{internal}
directions. At $n=1$ these are $e_2,e_3$ and $\wI=e_2e_3=J$.

\begin{theorem}[Selection of the imaginary unit]
\label{thm:i}
Inside $\End_\R(S_n)$ the joint commutant of the interior Lorentz algebra and of all internal bivectors is
$\operatorname{span}\{\id,\Omega,\wI,\Omega\wI\}\cong\C\otimes_\R\C\cong\C\oplus\C$, at every rung. Its solutions of
$x^2=-\id$ are the four isolated points $\pm\Omega$, $\pm\wI$. Of these only $\pm\wI$ commute with the interior gamma
matrices, from which the kinetic operator is built. Hence the imaginary unit is $i=\pm\wI$, and the chirality is
$\gamma_5=\mp\Omega\wI=\mp\Gamma$ with
\[
\Gamma:=\wI\,\Omega,\qquad \Gamma^2=\id .
\]
\end{theorem}

\begin{proof}
By \eqref{eq:bladecomm}, commuting with the six interior bivectors and with every internal bivector is a parity condition on
each blade, and the blades that satisfy all of them are exactly the four listed, at every rung. A four-dimensional
commutative real algebra with two commuting square roots of $-\id$ and no nilpotents is $\C\oplus\C$, whose square roots of
$-1$ are those four points. Finally $\Omega$ is the product of the four interior generators, so it anticommutes with each of
them, while $\wI$ contains no interior index and commutes with each.
\end{proof}

From now on $J$ denotes $\wI$; at one constituent this is $e_2e_3$, the generator of \cref{prop:6plus8plus1}. The selection
depends on its conditions, and it is worth saying exactly how.

\begin{remark}[The ambient algebra is part of the statement]
\label{rem:ambient}
If commutation with the internal bivectors is dropped, then at $n=1$ the commutant of the four gamma matrices in $M_8(\R)$ is the
quaternion algebra
\[
\HH=\operatorname{span}\{\id,\ J,\ e_2\Omega,\ e_3\Omega\},
\]
and its square roots of $-\id$ form a two-sphere. One further datum removes that freedom. The volume element of the temporal triple,
$T:=e_1e_2e_3$, depends only on the splitting $V=V^+\oplus V^-$ and on an orientation of $V^+$; it commutes with $J$ and anticommutes
with $e_2\Omega$ and $e_3\Omega$ by \eqref{eq:bladecomm}, so
\[
\{X\in\HH:\ [X,T]=0\}=\R\,\id\oplus\R\,J ,
\]
and $J=\gamma^0T$ is recovered from the clock and the temporal orientation, up to the sign of that orientation. Higher up the
tower the commutant of the gamma matrices alone acquires a matrix multiplicity algebra; at $n=3$ the Lorentz commutant is
$M_{128}(\C)$, whose own complex structure is $\Omega$ and not $\wI$ (\cref{prop:config}). Selection tracks the conditions imposed,
and every uniqueness statement in this paper is made with them.
\end{remark}

That the imaginary unit is not adjoined but read off the algebra is the cleanest structural result at one constituent. It has an
independent counterpart: on a real Hilbert space, Poincar\'e symmetry alone can force the complex structure \cite{MorettiOppio}.

% --- source: tex/s04_dirac#3 (Real Dirac theory is complex Dirac theory)
\subsection{Real Dirac theory is complex Dirac theory}

\begin{theorem}[Exactness of the realification]
\label{thm:realification}
The $\C$-linear algebra generated by the interior gamma matrices on $(S,J)\cong\C^4$ is all of $M_4(\C)$. Consequently the real
Maxwell--Dirac system on $\Cl(3,3)$ is ordinary complex Maxwell--Dirac theory written in real variables.
\end{theorem}

\begin{proof}
$J$ commutes with the four $\gamma^\mu$. The sixteen monomials in the $\gamma^\mu$ and their $J$-multiples are thirty-two
linearly independent real matrices commuting with $J$, and the commutant of $J$ in $M_8(\R)$ is $\End_\C(\C^4)=M_4(\C)$, of real
dimension thirty-two.
\end{proof}

\begin{corollary}
\label{cor:nopredict}
Anything derivable from the real formulation alone was derivable from the complex one, and conversely. The gauged system is quantum
electrodynamics in real variables, and no prediction can follow from the rewriting.
\end{corollary}

This is the deflationary principle of the introduction. It does not make the arena uninteresting; it says where its content can
lie. Everything below that depends only on the four gamma matrices and $J$ is inherited. What is not inherited is what the arena
adds: the selection of $J$, the surplus plane and its uses, the pairing, and the tower.

% --- source: tex/s04_dirac#4 (Admissible bilinears)
\subsection{Admissible bilinears}

A field $\psi$ with values in $S$ is real. If its components anticommute, a term $\psi^{\!\top}A\psi$ vanishes when $A$ is
symmetric, and $\psi^{\!\top}A\,\partial\psi=\tfrac12\partial(\psi^{\!\top}A\psi)$ is a total derivative when $A$ is antisymmetric. A
mass term therefore needs an antisymmetric kernel, and a first-order kinetic term a symmetric one. With the clock $\gamma^0=e_1$:

\begin{proposition}[Kernels of fixed symmetry]
\label{prop:kernels1}
The mass form $\gamma^0J$ is antisymmetric. The kinetic kernels
\[
K^\mu:=(\gamma^0J)(J\gamma^\mu)=-\gamma^0\gamma^\mu,\qquad K^0=-\id,
\]
are symmetric, while $(\gamma^0J)\gamma^\mu$ is antisymmetric. The kinetic form must therefore carry the factor $J$: its $J$-free analogue
defines a total derivative and constrains nothing a field equation sees.
\end{proposition}

The mass form $\gamma^0J=e_1e_2e_3=T$ is the temporal volume element of \cref{rem:ambient}. It is not the only admissible mass.

\begin{proposition}[Scalar and pseudoscalar mass]
\label{prop:massangle}
The antisymmetric elements of $M_8(\R)$ that commute with $J$ and whose bilinear form is invariant under the interior Lorentz algebra ($L^{\!\top}A+AL=0$ for every $L\in\lint$) form the
two-dimensional space $\operatorname{span}\{T,\ T\Omega\}$. Used as the kernel of the Dirac bilinear, the second contributes nothing
to the kinetic term, because $T\Omega J\gamma^\mu$ is antisymmetric for every $\mu$; the two differ only in the mass term. $T$ is a
scalar mass and $T\Omega$ a pseudoscalar one, since $\Omega^2=-\id$ and $\Omega$ anticommutes with every $\gamma^\mu$, so that $\Omega$
is the real form of $i\gamma_5$. The family $T(\cos\vartheta+\sin\vartheta\,\Omega)=Te^{\vartheta\Omega}$ is the chiral mass angle.
\end{proposition}

So $T$ is singled out among admissible masses by the requirement that the mass be scalar, and that requirement is an input. The
central element $\Omega$ will reappear in \cref{sec:discrete} as the axial $U(1)$.

\paragraph{The Hamiltonian.} The free Dirac Hamiltonian at momentum $p$ is
\begin{equation}
\label{eq:Hp}
H(p)=\gamma^0\gamma^ap_a+m\gamma^0,\qquad H(p)^2=E^2\id,\qquad E^2=|\bp|^2+m^2 ,
\end{equation}
with spectral projectors $\Lambda_\pm(p)=\tfrac12(\id\pm H(p)/E)$. Both $\gamma^0=e_1$ and $\gamma^0\gamma^a=e_1e_{a+3}$ are symmetric, so
$H(p)$ is a symmetric matrix that commutes with $J$; its eigenvalues are $\pm E$, each with real multiplicity four.

\begin{example}
\label{ex:Hp}
At $\bp=(0,0,3)$ and $m=4$, $H=3e_1e_6+4e_1$ satisfies $H^2=(9+16)\id=25\,\id$, because $e_1e_6$ and $e_1$ anticommute and
$(e_1e_6)^2=+\id$. The eigenvalues are $\pm5$; each eigenspace is a complex plane for $J$, and the helicity
$\Sigma_3=J\,\tfrac12e_4e_5$ commutes with $H$ and splits each of them into two real planes.
\end{example}

% --- source: tex/s04_dirac#5 (A first-order real form of free evolution)
\subsection{A first-order real form of free evolution}

\label{sec:firstorder}
For a commuting real field the same data give a first-order symplectic system. Take
\[
\cL_1=\tfrac12\Psi^{\!\top}J\dot\Psi-\tfrac12\Psi^{\!\top}H(p)\Psi .
\]
Its Euler--Lagrange equation is $\dot\Psi=-JH(p)\Psi$. Since $JH$ is antisymmetric with $(JH)^2=-E^2\id$, the flow lies in $SO(8)$ and
\[
\Psi(\lambda)=\cos(E\lambda)\,\Psi_0-\frac{\sin(E\lambda)}{E}\,JH\Psi_0 .
\]
Reading $J$ as $i$ turns this into $i\dot\psi=h\psi$ with $h=\alpha\cdot\bp+\beta m$, the free Dirac Hamiltonian at fixed momentum: an exact
embedding of free Dirac evolution in a real symplectic form, with eight second-class constraints reducing a sixteen-dimensional phase
space to eight. Three things norm conservation does not buy. The energy $\tfrac12\Psi^{\!\top}H\Psi$ is unbounded below, since $\Tr H=0$
and $H^2=E^2\id$ force the spectrum $\{\pm E\}$ with equal multiplicity. For anticommuting components the action collapses to a boundary
term, $\Psi^{\!\top}H\Psi=0$ because $H$ is symmetric and $\Psi^{\!\top}J\dot\Psi=\tfrac12\tfrac{d}{d\lambda}(\Psi^{\!\top}J\Psi)$ because $J$ is
antisymmetric, which is why a fermionic Lagrangian must use the kernels of \cref{prop:kernels1}. And the bilinear $\Psi^{\!\top}J\dot\Psi$
is not preserved by the boosts. The real form of Dirac theory in which a bivector commuting with the gamma matrices plays the role of $i$
is old \cite{Hestenes67}; what is specific here is that the bivector is the surplus rotation $J$, selected by \cref{thm:i}.

\section{The real Dirac Lagrangian}
\label{sec:w:therealdiraclagrangian}

% --- source: two_readings/r1_lagrangians#4 (The Dirac Lagrangian)
\subsubsection{The Dirac Lagrangian}
\label{r1:sec:dirac}

\cref{sec:dirac} fixes the kernels and the conditions they must meet, not a normalised Lagrangian: the mass form $T=\gamma^0J$
is antisymmetric, the kinetic kernels $K^\mu=(\gamma^0J)(J\gamma^\mu)=-\gamma^0\gamma^\mu$ are symmetric with $K^0=-\id$, and
$(\gamma^0J)\gamma^\mu$ is antisymmetric (\cref{prop:kernels1}; r1-14, which also checks $T^2=-\id$, $[K^\mu,J]=[T,J]=0$ and that
$K^\mu J$ is antisymmetric). The normalised one-constituent Lagrangian is that of \cref{mb:massless}(iii). For a real Grassmann
field $\psi:\R^{1,3}\to S\otimes\Lambda$, in Heaviside--Lorentz units with $\hbar=c=1$,
\begin{equation}
\label{r1:eq:LD}
L_D=\tfrac12\psi^{\!\top}\bigl(K^\mu D_\mu-mT\bigr)\psi=\tfrac12\psi^{\!\top}T\bigl(J\gamma^\mu D_\mu-m\bigr)\psi,\qquad
D_\mu=\partial_\mu+qA_\mu J ,
\end{equation}
the real form of $\bar\psi(i\gamma^\mu D_\mu-m)\psi$ with $J$ in the place of $i$ \cite{Hestenes67}, the charge $q$ being dimensionless in
these units ($D_\mu=\partial_\mu+\tfrac q\hbar A_\mu J$ with $[q]=\mathrm C$ in SI); in SI units the density is
$\tfrac12\hbar c\,\psi^{\!\top}(K^\mu\partial_\mu-\tfrac{mc}\hbar T)\psi$ with $x^0=ct$ and $[\psi]=\mathrm m^{-3/2}$ (r1-36).

\begin{proposition}[The real Dirac field]
\label{r1:prop:dirac}
\begin{enumerate}[label=(\roman*),leftmargin=2.2em]
\item For arbitrary real $8\times8$ kernels, the Euler--Lagrange expression of $\tfrac12\psi^{\!\top}A^\mu\partial_\mu\psi-\tfrac12\psi^{\!\top}M\psi$
is $A^\mu_{\rm sym}\partial_\mu\psi-M_{\rm anti}\psi$. A symmetric mass kernel gives a vanishing bilinear and an antisymmetric kinetic
kernel a total derivative.
\item The field equation of \eqref{r1:eq:LD} at $q=0$ is $K^\mu\partial_\mu\psi=mT\psi$, that is $T(J\gamma^\mu\partial_\mu-m)\psi=0$,
equivalently $J\gamma^\mu\partial_\mu\psi=m\psi$ or $\gamma^\mu\partial_\mu\psi=-mJ\psi$; since
$(\gamma^\mu\partial_\mu-mJ)(\gamma^\nu\partial_\nu+mJ)=\Box+m^2$, each component obeys $(\Box+m^2)\psi=0$.
\item $L_D$ is invariant under $\psi\mapsto e^{\theta J}\psi$; $D_\mu$ is covariant under $\psi\mapsto e^{-q\lambda J}\psi$,
$A\mapsto A+d\lambda$, since $e^{-\theta J}$ preserves every $K^\mu$ and $T$. The Noether current $j^\mu=\tfrac12\psi^{\!\top}K^\mu J\psi$
satisfies $\partial_\mu j^\mu=\psi^{\!\top}J\,\mathcal E_D$ identically, with $\mathcal E_D=K^\mu\partial_\mu\psi-mT\psi$; $j^0=-\tfrac12\psi^{\!\top}J\psi$.
\item The canonical tensor $\Theta^\mu{}_\nu=\tfrac12\psi^{\!\top}K^\mu\partial_\nu\psi-\delta^\mu_\nu L_D$ satisfies
$\partial_\mu\Theta^\mu{}_\nu=-(\partial_\nu\psi)^{\!\top}\mathcal E_D$ identically, and its energy density is
$\Theta^0{}_0=\tfrac12\psi^{\!\top}\bigl(\gamma^0\gamma^a\partial_a+mT\bigr)\psi$; on a mode with $\partial_a\psi=Jp_a\psi$ it is
$\tfrac12\psi^{\!\top}JH(p)\psi$, with $H(p)=\gamma^0\gamma^ap_a+m\gamma^0$ the free Hamiltonian of equation~\cref{eq:Hp}.
\item On shell $\Theta^{\mu\nu}-\Theta^{\nu\mu}=\partial_\rho S^{\rho\mu\nu}$ with the totally antisymmetric spin tensor
$S^{\rho\mu\nu}=\tfrac14\psi^{\!\top}\gamma^0\gamma^{[\rho}\gamma^\mu\gamma^{\nu]}\psi$, so the Belinfante tensor
$T_B^{\mu\nu}=\Theta^{\mu\nu}-\tfrac12\partial_\rho S^{\rho\mu\nu}=\tfrac14\psi^{\!\top}(K^\mu\partial^\nu+K^\nu\partial^\mu)\psi-\eta^{\mu\nu}L_D$ is
symmetric, conserved on shell, and has $T_B^{00}=\Theta^{00}$.
\end{enumerate}
\end{proposition}

\begin{proof}
For a Grassmann $\psi$, $\psi^{\!\top}M\delta\psi=\delta\psi^{\!\top}M\psi$ when $M$ is antisymmetric and $\psi^{\!\top}A\partial\delta\psi=
\partial(\psi^{\!\top}A\delta\psi)+\delta\psi^{\!\top}A^{\!\top}\partial\psi$; collecting the left variations gives (i), which the jet algebra
confirms for random integer kernels (r1-15). (ii) $K^\mu=TJ\gamma^\mu$ and $T$ is invertible (r1-16). (iii) $(K^\mu J)^{\!\top}=-K^\mu J$,
and $\partial_\mu j^\mu=\psi^{\!\top}K^\mu J\partial_\mu\psi=\psi^{\!\top}J(\mathcal E_D+mT\psi)$ with $\psi^{\!\top}JT\psi=0$ because $JT$ is
symmetric (r1-17). (iv) Exchanging Grassmann factors and using the symmetry of $K^\mu$;
$\gamma^0\gamma^aJ=J\gamma^0\gamma^a$ and $mT=Jm\gamma^0$ give the mode form (r1-18). (v) The coefficient $\tfrac14$ is solved
for, and the identity holds for all six pairs $(\mu,\nu)$ after the time derivatives are eliminated with the field equation (r1-19).
\end{proof}

Item (i) is fixed and turns the admissibility conditions of \cref{prop:kernels1} into a theorem of the variation: a Grassmann field
sees only the symmetric part of a first-order kinetic kernel and only the antisymmetric part of a mass kernel. Items (ii)--(iii) are
restated (\cref{sec:discrete}, where the Dirac operator is written $J\gamma^\mu\partial_\mu-m$); the sign of the mass term in
$\gamma^\mu\partial_\mu\psi=-mJ\psi$ follows from $K^0=-\id$ and $T=\gamma^0J$, and $m\mapsto-m$ passes to the opposite sign. Items
(iv)--(v) are inherited \cite{Belinfante40,WeinbergQFT}. Equation~\cref{eq:Saction}, on the rank-$64$ block of $P_*$ in the
three-constituent module, uses the kinetic kernels $P_*\gamma^0\gamma^\mu P_*$, which are minus the compressed $K^\mu$, with $K^0=+\id$ on its
block, and the Yukawa term $+\tfrac y2\psi^{\!\top}\Phi\psi$ in place of a mass term: with $y\Phi$ in the place of $mT$ its fermion Lagrangian is
minus \eqref{r1:eq:LD}, and its field equation $P_*\gamma^0\gamma^\mu P_*D_\mu\psi+y\Phi\psi=0$ has the same form (r1-16). For Grassmann
fields $\Theta^0{}_0$ is not a number, and the sign of the energy is decided only after quantisation. On the state space that
\crefrange{thm:module}{thm:spectra} construct from outside the real module, the one-momentum Hamiltonian has the spectrum
$\{0,\pm E,\pm2E\}$ with multiplicities $6,4,4,1,1$, symmetric about zero and bounded below (r1-37); a positive energy requires measuring
it from the lowest state, the usual choice of vacuum.

% --- source: tex/s05_discrete#0 ((head))


% --- source: tex/s05_discrete#1 (Discrete symmetries, the massless gauge algebra, and the sta)
\section{Discrete symmetries, the massless gauge algebra, and the states}
\label{sec:discrete}

% --- source: tex/s05_discrete#2 (Charge conjugation, parity and time reversal)
\subsection{Charge conjugation, parity and time reversal}

At one constituent write $U=\langle e_2,e_3\rangle$ for the surplus plane, $\Omega=\gamma^0\gamma^1\gamma^2\gamma^3$, and
$\Gamma=J\Omega$. An operator $X$ on $S$ is \emph{linear} when $XJ=JX$ and \emph{antilinear} when $XJ=-JX$. It intertwines the free
Dirac operator $J\gamma^\mu\partial_\mu-m$ when it commutes with every $J\gamma^\mu$. Parity is required to be linear and time
reversal antilinear, and each must satisfy $XH(p)X^{-1}=H(-p)$ for every $p$ and $m$.

\begin{theorem}[Charge conjugation is the surplus plane]
\label{thm:C}
The antilinear operators that intertwine the Dirac operator form exactly the two-dimensional space $U$. Every unit vector $C\in U$
satisfies
\[
CH(p)C^{-1}=-H(-p),\qquad CJC^{-1}=-J,\qquad C\Gamma C^{-1}=-\Gamma,\qquad C\Lambda_+(p)C^{-1}=\Lambda_-(-p),
\]
commutes with every interior Lorentz generator, and preserves the mass form and every kinetic kernel $K^\mu$. Hence $C$ is unique up
to a gauge rotation $e^{\theta J}$, and the operator known as the Majorana real structure is charge conjugation.
\end{theorem}

\begin{proof}
If $XJ=-JX$, then $X(J\gamma^\mu)=(J\gamma^\mu)X$ is equivalent to $X\gamma^\mu=-\gamma^\mu X$. The conditions are linear, and
conjugation by a blade sends every blade to plus or minus itself, so the solution space is spanned by the blades that anticommute
with $J$ and with $e_1,e_4,e_5,e_6$. By \eqref{eq:bladecomm} a blade $e_I$ anticommutes with a vector $e_a$ exactly when
$|I|-[a\in I]$ is odd. For odd $|I|$ this forces $I\subset\{2,3\}$, hence $I=\{2\}$ or $\{3\}$, and both anticommute with $J$. For
even $|I|$ it forces $\{1,4,5,6\}\subset I$, hence $e_I=\Omega$ or $e_1\cdots e_6$, and both commute with $J$. The solutions are
therefore $e_2$ and $e_3$. The listed identities follow from the same rule, since $e_2,e_3$ anticommute with every $\gamma^\mu$, with
$J$ and with $\Gamma$, and commute with every interior bivector.
\end{proof}

\begin{corollary}[Equal mass, equal spin, opposite charge]
$C$ commutes with every generator of the interior Lorentz algebra, so it preserves the Casimir $-\sum_aS_a^2=\tfrac34\id$, and
$(CH(p)C^{-1})^2=E^2\id$. The Majorana condition $\Psi=C\Psi$ makes the current vanish: a particle that is its own antiparticle is
neutral.
\end{corollary}

\begin{theorem}[Parity and time reversal]
\label{thm:PT}
The admissible parities form the space $\langle\gamma^0,\gamma^0J\rangle$, and $P^2=+\id$ holds only at $P=\pm\gamma^0$, while $P=\pm\gamma^0J$
gives $P^2=-\id$. The admissible time reversals form $U\cdot\gamma^1\gamma^2\gamma^3$, so $T=C\gamma^1\gamma^2\gamma^3$ with
\[
T^2=-\id .
\]
The sign is not a convention: $B=\gamma^1\gamma^2\gamma^3$ has $B^2=+\id$ and anticommutes with every vector $C$ of the clock plane, so
$T^2=CBCB=-C^2B^2=-\|C\|^2\id$ for every choice of clock in the positive triple.
\end{theorem}

\begin{theorem}[$CPT$ lies on a circle]
\label{thm:CPT}
The product of the three operators depends on the parity chosen:
\[
P=\gamma^0:\quad CPT=-\Omega,\ (CPT)^2=-\id;\qquad\qquad P=\gamma^0J:\quad CPT=+\Gamma,\ (CPT)^2=+\id .
\]
What is convention-free is the orbit of $\Omega$ under $U(1)_J$, $\theta\mapsto\cos\theta\,\Omega+\sin\theta\,J\Omega$, whose four real points
$\pm\Omega,\pm\Gamma$ are the spacetime pseudoscalar and the chirality.
\end{theorem}

\begin{figure}[t]
\centering
\begin{tikzpicture}[scale=1.25,>=Stealth]
\draw[thick,gray] (0,0) circle (1.3);
\foreach \a/\l/\c in {0/{+\Omega}/scol,90/{+\Gamma=J\Omega}/tcol,180/{-\Omega}/scol,270/{-\Gamma}/tcol}{
  \fill[\c] (\a:1.3) circle (2pt);
  \node[\c] at (\a:1.85) {$\l$};
}
\draw[->,thick] (20:0.75) arc (20:70:0.75);
\node at (45:0.42) {\small $e^{\theta J}$};
\node[align=left,anchor=west] at (2.7,0.55) {\small $P=\gamma^0$: \quad $CPT=-\Omega$,\ $(CPT)^2=-\id$};
\node[align=left,anchor=west] at (2.7,-0.05) {\small $P=\gamma^0J$: \ $CPT=+\Gamma$,\ $(CPT)^2=+\id$};
\node[align=left,anchor=west] at (2.7,-0.65) {\small both lie on the $U(1)_J$ orbit of $\Omega$};
\end{tikzpicture}
\caption{The $CPT$ circle at one constituent (\cref{thm:CPT}). The product of the three discrete operators is the spacetime
pseudoscalar or the chirality according to the parity chosen; the two lie on one orbit of the gauge rotation $e^{\theta J}$.}
\label{fig:cpt}
\end{figure}

These are statements about explicitly constructed operators of the free theory; no $CPT$ theorem is proved, and the interacting
theory is not constructed.

% --- source: tex/s05_discrete#3 (The massless gauge algebra at one constituent)
\subsection{The massless gauge algebra at one constituent}

A gauge generator must leave the kinetic form invariant. Solving that condition over all of $\End_\R(S)$, with no ansatz, gives the
following.

\begin{proposition}[Gauge algebra from the kinetic forms]
\label{prop:gauge1}
The solutions of $X^{\!\top}K^\mu+K^\mu X=0$ for $\mu=0,\dots,3$ form
\[
\mathfrak g_{\rm massless}=\R\,\Omega\ \oplus\ \operatorname{span}\{J,\ e_2\Omega,\ e_3\Omega\}\ \cong\ \uu1_\Omega\oplus\su2 .
\]
The condition at $\mu=0$ is $X^{\!\top}=-X$, and the conditions at $\mu=a$ then say that $X$ commutes with the three boosts
$\gamma^0\gamma^a$ and hence with the whole interior Lorentz algebra. The trace form is $-8\,\id$ on this basis, so the algebra is
compact. Adding invariance of the mass form, $X^{\!\top}(\gamma^0J)+(\gamma^0J)X=0$, leaves exactly $\R J$. The algebra meets the span of the
pole operators in zero.
\end{proposition}

The $\su2$ factor contains the fermion-number direction $J$ and two generators that anticommute with $J$; it therefore mixes a field with
its conjugate and is the Pauli--G\"ursey algebra of a massless Dirac field \cite{Pauli57,Gursey58}. It is a symmetry of the free massless
theory, not a gauge candidate on this matter: its two Weyl components form a single doublet, and gauging $SU(2)$ on an odd number of
doublets is forbidden by the global anomaly \cite{Witten82}. The central factor is the axial one. Since $\Gamma=J\Omega$, $\Omega=-J\Gamma$,
which with $i=J$ and the chirality $\gamma_5=-\Gamma$ of \cref{thm:i} is the real form of $i\gamma_5$, as in \cref{prop:massangle}; $U(1)_\Omega$ is broken by a Dirac mass and is anomalous on this matter content \cite{Bilal,AlvarezGaumeWitten}.
The electromagnetic $\uu1$ is therefore singled out twice at one constituent: as the whole internal centraliser
(\cref{prop:6plus8plus1}) and as what survives of the massless gauge algebra once a mass is admitted.

% --- source: tex/s05_discrete#4 (The state space)
\subsection{The state space}

\begin{theorem}[The real module is not a state space]
\label{thm:module}
Eight $16\times16$ integer matrices $\theta_i$, antisymmetric with $\theta_i^2=-\id$ and pairwise anticommuting, realise the canonical
anticommutation relations of the eight real field components exactly. Their commutant in $M_{16}(\R)$ is $\R\,\id$, because the $256$
Clifford monomials span the whole matrix algebra. Hence no internal complex structure commutes with all eight field operators, and the
quadratic observables built from them are real antisymmetric matrices with imaginary spectra. The quantum state space has to be formed
from outside, as $\mathcal F=\C\otimes_\R\R^{16}$, with charge $Q=-i\rho(J)$ and Hamiltonian $H=-i\rho(JH(p))$, where
$\rho(A)=-\tfrac14\sum_{ij}A_{ij}\theta_i\theta_j$ lifts an antisymmetric $8\times8$ matrix to the module.
\end{theorem}

It is the complex structure, not the reality, that is missing, which is the lesson of \cref{rem:ambient} arriving one level up.

\begin{theorem}[The free spectrum of one momentum]
\label{thm:spectra}
The lifted generators have characteristic polynomials
\[
\chi_{2\rho(J)}(x)=x^6(x^2+4)^4(x^2+16),\qquad \chi_{2\rho(JH)}(x)=x^6(x^2+4E^2)^4(x^2+16E^2),
\]
so $\spec Q=\{0,\pm1,\pm2\}$ and $\spec H=\{0,\pm E,\pm2E\}$, with multiplicities $6,4,4,1,1$. The joint eigenspaces of $(H,Q)$ are
\begin{center}
\begin{tabular}{@{}lll@{}}
\toprule
$H$ & dimension & charges\\
\midrule
$-2E$ & $1$ & $q=0$\\
$-E$ & $4$ & $q=\pm1$, twice each\\
$0$ & $6$ & $q=\pm2$ once each, $q=0$ four times\\
$+E$ & $4$ & $q=\pm1$, twice each\\
$+2E$ & $1$ & $q=0$\\
\bottomrule
\end{tabular}
\end{center}
so $H+2E$ has spectrum $0,E,2E,3E,4E$ with multiplicities $1,4,6,4,1$.
\end{theorem}

\begin{figure}[t]
\centering
\begin{tikzpicture}[xscale=1.2,yscale=0.95]
\foreach \y/\l in {-2/{-2E},-1/{-E},0/{0},1/{+E},2/{+2E}}{
  \draw[gray!60] (-2.6,\y) -- (2.6,\y);
  \node[left] at (-2.7,\y) {\small $\l$};
}
\node at (0,2.75) {\small charge $q$};
\foreach \q in {-2,-1,0,1,2}{\node at (\q,-2.7) {\small $\q$};}
% ground state
\fill[tcol] (0,-2) circle (2.4pt); \node[above right] at (0,-2) {\scriptsize 1};
% -E level
\fill[tcol] (-1,-1) circle (2.4pt); \fill[tcol] (1,-1) circle (2.4pt);
\node[above] at (-1,-1) {\scriptsize 2}; \node[above] at (1,-1) {\scriptsize 2};
% 0 level
\fill[tcol] (-2,0) circle (2.4pt); \fill[tcol] (2,0) circle (2.4pt); \fill[tcol] (0,0) circle (2.4pt);
\node[above] at (-2,0) {\scriptsize 1}; \node[above] at (2,0) {\scriptsize 1}; \node[above] at (0,0) {\scriptsize 4};
% +E
\fill[tcol] (-1,1) circle (2.4pt); \fill[tcol] (1,1) circle (2.4pt);
\node[above] at (-1,1) {\scriptsize 2}; \node[above] at (1,1) {\scriptsize 2};
% +2E
\fill[tcol] (0,2) circle (2.4pt); \node[above right] at (0,2) {\scriptsize 1};
\node[align=left,anchor=west] at (3.0,-2) {\small vacuum: unique, $q=0$};
\node[align=left,anchor=west] at (3.0,-1) {\small one particle or antiparticle};
\node[align=left,anchor=west] at (3.0,0) {\small two excitations};
\node[align=left,anchor=west] at (3.0,1) {\small three excitations};
\node[align=left,anchor=west] at (3.0,2) {\small filled};
\end{tikzpicture}
\caption{The Fock space of one momentum (\cref{thm:spectra}). Each dot is a joint eigenspace of energy and charge, labelled by its
dimension. The multiplicities $1,4,6,4,1$ are those of four fermionic modes: two particle and two antiparticle states of one Dirac
flavour, read off rather than inferred.}
\label{fig:fock}
\end{figure}

The multiplicities $1,4,6,4,1$ are the binomial coefficients of a four-mode Fock space, as a single Dirac flavour with two particle and
two antiparticle states must give. The ground state at $-2E$ is the zero-point energy $-\tfrac12\sum E$ of the four modes; summed over
momenta it is the familiar fermion term of the one-loop effective potential. A projector built from $\rho(JH)^2$ could not have
produced the table: squaring loses the sign of the energy and conflates the vacuum with the filled state.

\begin{proposition}[Conjugation on states]
\label{prop:Cstates}
$C\in SO(8)$ lifts to $\mathrm{Spin}(8)$ as a product of four integral reflections. The lift is unitary on $\mathcal F$, anticommutes with
$Q$ and carries $H(p)$ to $H(-p)$, so $\langle\psi|Q|\psi\rangle=0$ for every $C$-eigenstate. The stronger statement $Q|\psi\rangle=0$ is
false: the $+1$ eigenspace of $C$ has real dimension eight and $Q$ annihilates only a three-dimensional part of it.
\end{proposition}

\begin{remark}[Many-body chirality is a group lift]
$\rho$ is a homomorphism of Lie algebras only, so the product $\rho(J)\rho(\Omega)$, of rank eight, is not a $\mathrm{Spin}(8)$ element and
is not the chirality of the module. The correct object is $\Gamma_{\mathcal F}=\exp[\tfrac\pi2(\rho(J)+\rho(\Omega))]$, an involution of
$\mathcal F$ whose eigenvalue on a joint eigenvector is $i^{q_J+q_\Omega}$; the sectors are labelled by the parity of
$\tfrac12(q_J+q_\Omega)$.
\end{remark}

